\section{Optimization}

\subsection{Caching}

            Faster CPU's does not mean faster execution. The CPU might compute almost instantly but the loading and
            storing of the exectution data/result forms a bottleneck, as this is one of the slowest operations.

            \begin{definition}

\textbf{Cache} is a fast memory that preloads and stores frequently used data form the memory to reduce
                latency.
            
\end{definition}

            The usage of cache reduces stalls (times in which the CPU can not execute as the data is not yet present).

            \subsection{Vectorization}

            With vectorized data we can use SIMD to exectute the same operation on multiple data points all at once.

            \begin{definition}

\textbf{SIMD} (Single Instruction Multiple Data) is a type of parallel processing where the same
                instruction
                is
                applied to multiple data points simultaneously by using multiple ALUs inside the CPU. The same
                instruction gets broadcasted to multiple ALU which execute in parallel.
            
\end{definition}

\subsection{Superscalar and Speculative Execution (ILP)}

\begin{definition}

\textbf{Superscalar execution} refers to out of order execution of independent instructions in parallel
                using
                multiple ALUs.
            
\end{definition}

\begin{definition}

\textbf{Speculative execution} refers to the execution of instructions before the actual need for them is
                known.
            
\end{definition}

\subsection{Pipelining}

            Pipelining increases throughput by overlapping the execution of multiple instructions in diffrent stages.

            \begin{example}

                As soon as I am finished washing my laundry and start hanging it, my neighbour can already use the
                machine and start wahsing his clothes. My neighbour does not have to wait for me to wash, hang, dry,
                fold, store my clothes.
            
\end{example}

\begin{definition}

\textbf{Throughput } is the amout of work that can be done by a system in a given period of time. It can
                be bounded by
                $$tp = \frac{1}{max(computationTime(stages))}$$
            
\end{definition}

\begin{definition}

\textbf{Latency } is the time needed to perform a given computation. It can be bound by
                $$L = numStages * max(computationTime(stages))$$
            
\end{definition}

\begin{definition}

                A pipeline is \textbf{balanced } if all stages take the same amount of time.
            
\end{definition}

            Balancing a pipeline can increase throughput.

            \begin{example}

                Assume washing is the stage that takes the longest. Having two small washing machines instead of a large
                one can increase through put as now my neighbours neighbour can start wahsing as soon as i am finished.
                This balances the pipeline.
            
\end{example}

\subsection{Parallel Performance}

\begin{definition}

\textbf{Scalability } describes how performance improves with additional resources (cores, memory, etc.).
            
\end{definition}

\begin{definition}

\textbf{Speedup } is the ratio of the execution time of the sequential algorithm to the execution time of
                the parallel algorithm.
                $$S_p = \frac{T_1}{T_p}$$
            
\end{definition}

\begin{definition}

\textbf{Efficiency } is the ratio of speedup to the number of processors.
                $$E_p = \frac{S_p}{p}$$
            
\end{definition}

\begin{remark}

                Speedup depends on the definition of $T_1$, is it the execution time of the optimized sequential
                algorithm or the unoptimized one?
            
\end{remark}

\begin{theorem}

                Amdahl's Law sets a limit on speedup, as the sequential part of the algorithm can not be parallelized.
                $$S(p) = \frac{1}{1-f + \frac{f}{p}}$$
            
\end{theorem}

\begin{theorem}

                Gustafson's Law proposes an optimistic view compared to Amdahl. It shows that we might not be able to
                shrink execution time but we can perform more operations in that same execution time.
                $$S_p = f + P(1 - f)$$
            
\end{theorem}

